\par\noindent\textbf{(1)} 静电均匀场中介电球有效极化率α=4πεm a³(εp−εm)/(εp+2εm)，忽略动态辐射修正。若诱导偶极p=αE且α为实数，周期平均势能U=−α|E0|²/4，因I=(1/2)nm c ε0|E0|²，故U∝−α I。
\par\noindent\textbf{(2)} 梯度力F\_grad=−∇U=(α/4)∇|E0|²，对Gaussian强度峰，α>0的球受向强场中心的横向回复力；α<0则被排向暗区。由于强度沿束轴也变化，三维稳定还取决于数值孔径和轴向散射压力，不能仅看横向势阱。
\par\noindent\textbf{(3)} 在Rayleigh极限散射截面σ\_sca=k\_m⁴|α|²/(6πεm²)，辐射压力数量级F\_sca≈nm Iσ\_sca/c沿传播方向。因α∝a³，梯度力对固定强度梯度∝a³，而散射力∝a⁶，所以小粒子更利于梯度捕获；两者比还含束腰长度和折射率对比。
\par\noindent\textbf{(4)} εp>εm时静态α>0，可在高数值孔径聚焦中心获得稳定光阱；εp<εm时倾向低强度环或暗阱。若a不再远小于λ，偶极模型不准，应改用Mie散射或电磁Maxwell应力张量计算总力。
